\subsection{Dispersed additive-update representation}
\label{sec:w1-dispersed}

\begin{definition}[$r$-round dispersed update stream]
\label{def:w1-dispersed-stream}
The cyclic specialization of \cref{def:phase-family} with modulus
$Q\geq3$, $K=Q-1$, $\Delta=2\pi/Q$, $r\geq2$ rounds, and
$\Sigma=\{+1\}$: for $x\in[K]_0^m$ the shares
$a^{(1)},\ldots,a^{(r)}\in\mathbb Z_Q^m$ satisfy
\begin{equation}
 \sum_{t=1}^r a_j^{(t)}=x_j\pmod Q
 \qquad(j\in[m]),
 \label{eq:w1-share-sum}
\end{equation}
the flat input is the round-major record sequence
$\mathsf{Update}(t,j,a_j^{(t)})$, and the semantic input is the addressed
aggregate vector $x$.
\end{definition}

\begin{proposition}[Semantic equivalence and token privacy]
\label{prop:w1-stream-equivalence}
Every valid flat stream of \cref{def:w1-dispersed-stream} implements the
same $U_x$ as the semantic input, up to global phase.  Moreover, a single
flat token never determines $x_j$: for every fixed $(t,j,a)$ and every
$b\in[K]_0$, there is a valid stream with $a_j^{(t)}=a$ and $x_j=b$.
Under the canonical sharing distribution in which the first $r-1$ shares
are independent uniform residues and the last is fixed by
\eqref{eq:w1-share-sum}, every individual residue is uniform on
$\mathbb Z_Q$ and independent of $x_j$.
\end{proposition}

\begin{proof}
For integer representatives there are $k_j\in\mathbb Z$ with
$\sum_ta_j^{(t)}=x_j+Qk_j$.  Commutativity and $Q\Delta=2\pi$ give
\begin{equation*}
 \prod_tR_{P_j}(\Delta a_j^{(t)})
 =R_{P_j}(\Delta x_j+2\pi k_j)
 =(-1)^{k_j}R_{P_j}(\Delta x_j).
\end{equation*}
The product of the signs is global.  For privacy, fix the named token and
choose any one of the other $r-1$ residues to enforce the desired sum;
the remaining residues are arbitrary.  The distributional claim follows
because a sum containing at least one independent uniform residue is
uniform, including the last share after conditioning on $x_j$.
\end{proof}

\begin{definition}[Forward-pass cut accounting]
\label{def:w1-pass-accounting}
Partition the round-major stream into Alice's prefix $A$ (rounds
$1,\ldots,r-1$) and Bob's suffix $B$ (round $r$).  A forward $p$-pass
compiler scans $AB$ from left to right at most $p=p(m)$ times.  There are
$p$ crossings $A\to B$ and at most $p-1$ rewind crossings $B\to A$.
At every crossing, the complete restart state is the five-field
serialization of \eqref{eq:info-budget}; rereadable input and every live
IR object are therefore included.  For a run on $I$, let $A_p(I)$ be the
number of bits
irreversibly committed to the final output while the machine scans $A$,
summed over all passes, and define
\begin{equation}
 \Sigma_p(I):=\sum_{t=1}^{p}B_{\mathrm{cut}}(c_{A\to B,t})
       +\sum_{t=1}^{p-1}B_{\mathrm{cut}}(c_{B\to A,t}).
 \label{eq:w1-transcript-budget}
\end{equation}
Both are literal, self-delimiting transcript lengths.  Unbarred resource
symbols take maxima over valid inputs.  For randomized compilers, write
\begin{equation*}
 \overline A_p:=\max_I\mathbb E_\rho A_p(I,\rho),\qquad
 \overline\Sigma_p:=\max_I\mathbb E_\rho\Sigma_p(I,\rho),
\end{equation*}
and let
$S:=\max_I\mathbb E_\rho[\max_cB_{\mathrm{cut}}(c;I,\rho)]$ over the
$2p-1$ named crossings.  Then
$\overline\Sigma_p\leq(2p-1)S$; the deterministic version is identical
without expectations.
\end{definition}

\begin{theorem}[One-pass zero-error mask bound]
\label{thm:w1-one-pass-mask}
Let a deterministic one-pass compiler be correct to
$\epsilon<\sin(\pi/(2Q))$ on every valid stream of
\cref{def:w1-dispersed-stream}.  Then
\begin{equation}
 A_1+B_{\mathrm{cut}}(c_{A\to B,1})
 \geq\left\lceil m\log_2Q\right\rceil.
 \label{eq:w1-one-pass-mask}
\end{equation}
\end{theorem}

\begin{proof}
For each $u\in\mathbb Z_Q^m$, choose a canonical prefix whose aggregate
is $u$ (put $u$ in its first round and zero in the other prefix rounds).
The committed prefix and restart state after this prefix must distinguish
all $Q^m$ choices.  Indeed, for $u\neq u'$, choose one common final share
$v$ coordinatewise.  At each coordinate at most two residues of $v_j$
would make either $u_j+v_j$ or $u'_j+v_j$ equal to the single excluded
aggregate $Q-1$; since $Q\geq3$, an allowed $v_j$ exists.  Both completions
are valid, and they differ in every coordinate in which $u$ and $u'$ differ.
By \cref{thm:w1-exact-separation}, one circuit cannot be within
$\epsilon$ of both targets.  The messages are prefix-free under the
declared serializers, so Kraft's inequality gives
\eqref{eq:w1-one-pass-mask}.
\end{proof}
